\textbf{Phần 1. Trắc nghiệm lựa chọn}
\setcounterpageref{ex}{0}

\begin{ex} %[3H5N30]
Mi-li-mét viết tắt là
\choice
{$\text{m}$}
{\True $\text{mm}$}
{$\text{cm}$}
{$\text{dm}$}
\loigiai{
Mi-li-mét viết tắt là $\text{mm}$.\\
\textbf{Đáp án của câu hỏi là B}
}
\end{ex}

\begin{ex} %[3H5N30]
$1\text{ cm}$ bằng bao nhiêu mi-li-mét?
\choice
{$1\text{ mm}$}
{$100\text{ mm}$}
{$1000\text{ mm}$}
{\True $10\text{ mm}$}
\loigiai{
$1\text{ cm} = 10\text{ mm}$.\\
\textbf{Đáp án của câu hỏi là D}
}
\end{ex}

\begin{ex} %[3H5H30]
$1\text{ m}$ bằng bao nhiêu mi-li-mét?
\choice
{\True $1000\text{ mm}$}
{$10\text{ mm}$}
{$100\text{ mm}$}
{$10000\text{ mm}$}
\loigiai{
$1\text{ m} = 100\text{ cm} = 1000\text{ mm}$.\\
\textbf{Đáp án của câu hỏi là A}
}
\end{ex}

\begin{ex} %[3H5H30]
Đổi: $4\text{ cm } 5\text{ mm} = \dots\text{ mm}$.
\choice
{$405\text{ mm}$}
{$450\text{ mm}$}
{\True $45\text{ mm}$}
{$54\text{ mm}$}
\loigiai{
$4\text{ cm } 5\text{ mm} = 40 + 5 = 45\text{ mm}$.\\
\textbf{Đáp án của câu hỏi là C}
}
\end{ex}

\begin{ex} %[3H5H30]
Tính: $35\text{ mm} + 28\text{ mm} = \dots$
\choice
{$53\text{ mm}$}
{$65\text{ mm}$}
{$73\text{ mm}$}
{\True $63\text{ mm}$}
\loigiai{
$35 + 28 = 63\text{ mm}$.\\
\textbf{Đáp án của câu hỏi là D}
}
\end{ex}

\begin{ex} %[3H5H30]
Tính: $80\text{ mm} - 35\text{ mm} = \dots$
\choice
{\True $45\text{ mm}$}
{$40\text{ mm}$}
{$50\text{ mm}$}
{$55\text{ mm}$}
\loigiai{
$80 - 35 = 45\text{ mm}$.\\
\textbf{Đáp án của câu hỏi là A}
}
\end{ex}

\begin{ex} %[3H5V30]
Bề dày cuốn sách Toán 3 dày khoảng $8\text{ mm}$. Bề dày của 5 cuốn sách như thế xếp chồng lên nhau là
\choice
{$35\text{ mm}$}
{\True $40\text{ mm}$}
{$45\text{ mm}$}
{$50\text{ mm}$}
\loigiai{
$8 \times 5 = 40\text{ mm}$.\\
\textbf{Đáp án của câu hỏi là B}
}
\end{ex}

\begin{ex} %[3H5V30]
Một cái đinh dài $5\text{ cm}$. Người ta đóng ngập vào gỗ một đoạn $35\text{ mm}$. Phần đinh còn nhô ra ngoài dài là
\choice
{$10\text{ mm}$}
{$20\text{ mm}$}
{$25\text{ mm}$}
{\True $15\text{ mm}$}
\loigiai{
$5\text{ cm} = 50\text{ mm}$.\\
Phần nhô ra: $50 - 35 = 15\text{ mm}$.\\
\textbf{Đáp án của câu hỏi là D}
}
\end{ex}

\begin{ex} %[3H5V30]
Một con kiến bò từ điểm $A$ đến điểm $B$ dài $12\text{ cm}$, rồi bò tiếp đoạn $BC$ dài $35\text{ mm}$. Tổng quãng đường kiến bò là
\choice
{$145\text{ mm}$}
{$165\text{ mm}$}
{\True $155\text{ mm}$}
{$47\text{ mm}$}
\loigiai{
$12\text{ cm} = 120\text{ mm}$.\\
Tổng quãng đường: $120 + 35 = 155\text{ mm}$.\\
\textbf{Đáp án của câu hỏi là C}
}
\end{ex}

\begin{ex} %[3H5V30]
Đổi: $70\text{ mm} = \dots\text{ cm}$.
\choice
{\True $7\text{ cm}$}
{$700\text{ cm}$}
{$70\text{ cm}$}
{$0{,}7\text{ cm}$}
\loigiai{
$70 : 10 = 7\text{ cm}$.\\
\textbf{Đáp án của câu hỏi là A}
}
\end{ex}

\begin{ex} %[3H5V30]
Tính: $16\text{ mm} \times 4 = \dots$
\choice
{$54\text{ mm}$}
{\True $64\text{ mm}$}
{$74\text{ mm}$}
{$60\text{ mm}$}
\loigiai{
$16 \times 4 = 64\text{ mm}$.\\
\textbf{Đáp án của câu hỏi là B}
}
\end{ex}

\begin{ex} %[3H5C30]
Một thanh gỗ dài $1\text{ m}$. Bác thợ cắt thành các đoạn dài $200\text{ mm}$. Bác cắt được số đoạn gỗ là
\choice
{$4$ đoạn}
{$6$ đoạn}
{\True $5$ đoạn}
{$10$ đoạn}
\loigiai{
$1\text{ m} = 1000\text{ mm}$.\\
Số đoạn: $1000 : 200 = 5$ đoạn.\\
\textbf{Đáp án của câu hỏi là C}
}
\end{ex}

\textbf{Phần 2. Trắc nghiệm đúng sai}
\setcounterpageref{ex}{0}

\begin{ex} %[3H5H30]
Xét tính đúng sai của việc đổi đơn vị:
\choiceTF
{\True $1\text{ cm} = 10\text{ mm}$.}
{\True $5\text{ cm } 3\text{ mm} = 53\text{ mm}$.}
{$1\text{ m} = 100\text{ mm}$.}
{\True $1\text{ m} = 1000\text{ mm}$.}
\loigiai{
\begin{itemize}
    \item $1\text{ cm} = 10\text{ mm}$: Đúng.
    \item $5\text{ cm } 3\text{ mm} = 53\text{ mm}$: Đúng.
    \item $1\text{ m} = 1000\text{ mm} \neq 100\text{ mm}$. Ý 3 Sai, Ý 4 Đúng.
\end{itemize}
\begin{itemize}
    \item a) \textbf{Đúng}.
    \item b) \textbf{Đúng}.
    \item c) \textbf{Sai}.
    \item d) \textbf{Đúng}.
\end{itemize}
}
\end{ex}

\begin{ex} %[3H5V30]
Đoạn bút chì lúc mới mua dài $18\text{ cm}$, sau một thời gian gọt dùng còn lại $95\text{ mm}$. Xét tính đúng sai:
\choiceTF
{\True Độ dài bút chì ban đầu là $180\text{ mm}$.}
{\True Đoạn bút chì đã bị mòn gọt đi là $85\text{ mm}$.}
{Đoạn bút chì còn lại dài hơn đoạn đã gọt là $15\text{ mm}$.}
{\True Đoạn bút chì còn lại dài hơn đoạn đã gọt là $10\text{ mm}$.}
\loigiai{
\begin{itemize}
    \item $18\text{ cm} = 180\text{ mm}$. Đúng.
    \item Đã gọt: $180 - 95 = 85\text{ mm}$. Đúng.
    \item Hiệu: $95 - 85 = 10\text{ mm} \neq 15\text{ mm}$. Ý 3 Sai, Ý 4 Đúng.
\end{itemize}
\begin{itemize}
    \item a) \textbf{Đúng}.
    \item b) \textbf{Đúng}.
    \item c) \textbf{Sai}.
    \item d) \textbf{Đúng}.
\end{itemize}
}
\end{ex}

\textbf{Phần 3. Trả lời ngắn}
\setcounterpageref{ex}{0}

\begin{ex} %[3H5H30]
Đổi $6\text{ cm } 8\text{ mm}$ ra mi-li-mét.
\shortans{68}
\loigiai{
\textbf{Đáp số:} $68$\\
$60 + 8 = 68\text{ mm}$.
}
\end{ex}

\begin{ex} %[3H5H30]
Tính: $45\text{ mm} + 35\text{ mm}$.
\shortans{80}
\loigiai{
\textbf{Đáp số:} $80$\\
$45 + 35 = 80\text{ mm}$.
}
\end{ex}

\begin{ex} %[3H5V30]
Tính: $90\text{ mm} : 6$.
\shortans{15}
\loigiai{
\textbf{Đáp số:} $15$\\
$90 : 6 = 15\text{ mm}$.
}
\end{ex}

\begin{ex} %[3H5V30]
Một sợi dây đồng dài $10\text{ cm}$. Người ta cắt lấy $45\text{ mm}$. Hỏi đoạn dây còn lại dài bao nhiêu mi-li-mét?
\shortans{55}
\loigiai{
\textbf{Đáp số:} $55$\\
$100 - 45 = 55\text{ mm}$.
}
\end{ex}

\begin{ex} %[3H5V30]
Tính chu vi tam giác có các cạnh dài $20\text{ mm}, 30\text{ mm}, 40\text{ mm}$ theo mi-li-mét.
\shortans{90}
\loigiai{
\textbf{Đáp số:} $90$\\
$20 + 30 + 40 = 90\text{ mm}$.
}
\end{ex}

\begin{ex} %[3H5C30]
Một chồng gồm 8 quyển vở cùng loại dày tất cả $48\text{ mm}$. Hỏi mỗi quyển vở dày bao nhiêu mi-li-mét?
\shortans{6}
\loigiai{
\textbf{Đáp số:} $6$\\
$48 : 8 = 6\text{ mm}$.
}
\end{ex}
